\documentclass[11pt,a4paper]{article}
\usepackage[margin=25mm,headheight=14pt]{geometry}
\usepackage{fontspec}
\setmainfont{lmroman10-regular.otf}[BoldFont=lmroman10-bold.otf,ItalicFont=lmroman10-italic.otf,BoldItalicFont=lmroman10-bolditalic.otf]
\setsansfont{lmsans10-regular.otf}
\setmonofont{lmmono10-regular.otf}[Scale=MatchLowercase]
\newfontfamily\authorsymbol{Noto Sans Symbols}[Scale=MatchUppercase]
\usepackage{amsmath,amssymb,amsthm,mathtools,booktabs,array}
\usepackage{microtype}
\usepackage[hidelinks,unicode]{hyperref}
\usepackage{fancyhdr}
\pagestyle{fancy}\fancyhf{}
\fancyhead[L]{\small\itshape Exact state complexity of two-stage addition}
\fancyhead[R]{\small\thepage}\renewcommand{\headrulewidth}{0pt}
\setlength{\parindent}{1.2em}\setlength{\parskip}{0pt}
\setlength{\emergencystretch}{2em}
\newtheorem{theorem}{Theorem}
\newtheorem{proposition}[theorem]{Proposition}
\newtheorem{lemma}[theorem]{Lemma}
\newtheorem{corollary}[theorem]{Corollary}
\theoremstyle{definition}\newtheorem{definition}[theorem]{Definition}
\newcommand{\res}[1]{[#1]_B}
\newcommand{\pairs}{\{0,1\}^2}
\newcommand{\Q}{\mathcal Q_B}
\DeclareMathOperator{\im}{im}
\hypersetup{pdftitle={Exact State Complexity of Two-Stage Addition},pdfauthor={lmtlssss🜐},pdfsubject={Mathematics: finite functions and modular arithmetic},pdfkeywords={carry propagation, finite functions, state complexity, exact encoding}}
\title{\LARGE Exact State Complexity of Two-Stage Addition}
\author{lmtlssss{\authorsymbol 🜐}}
\date{27 October 2026}
\begin{document}
\maketitle\thispagestyle{plain}
\begin{abstract}
Three fixed radix-$B$ operands determine a function of two incoming carry
bits through two successive additions. We classify this function when its
outputs include the residue and both outgoing carries. For every integer
$B\geq2$, exactly $5B-4$ functions occur. An explicit normal form gives a
contiguous state code and a representative operand triple for each state.
We determine each class's multiplicity and the distribution of carry-image
sizes under uniform operands. The minimum fixed-length description is
$\lceil\log_2(5B-4)\rceil$ bits. In particular, three 32-bit operands admit
an exact 35-bit response code. Independently recoverable digit stages obey
a separate product formula; four byte stages require 42 bits. The proofs
are accompanied by complete-function enumeration and executable encoders.
\end{abstract}
\noindent\textit{Keywords.} Modular arithmetic; carry propagation; finite functions;
exact encoding; enumerative combinatorics.\par
\noindent\textit{DOI.} \href{https://doi.org/10.5281/zenodo.23698617}{10.5281/zenodo.23698617}.

\section{The response contract}
Fix an integer $B\geq2$. Write $\res z$ for the least nonnegative residue
of $z$ modulo $B$. The operands $C,S,H$ belong to $\{0,\ldots,B-1\}$.
They remain fixed while the incoming carry bits $(a,b)$ range over $\pairs$.
Define
\begin{align}
 u&=C+S+a,& c_1&=\lfloor u/B\rfloor,\label{eq:first}\\
 z&=\res u+H+b,& c_2&=\lfloor z/B\rfloor,\qquad \rho=\res z.\label{eq:second}
\end{align}
Both $c_1$ and $c_2$ are bits. The \emph{complete response} is the function
\[
 Q_{C,S,H}:\pairs\longrightarrow\{0,\ldots,B-1\}\times\pairs,
 \qquad (a,b)\longmapsto(\rho,c_1,c_2).
\]
Two operand triples are equivalent when their complete responses agree
at all four inputs. Let $\Q$ denote the set of these functions.
A state code identifies an element of $\Q$; it need not recover the
original operands.

Carry composition is an established method in arithmetic circuits.
Parallel-prefix addition uses it to combine local carry information
\cite{brentkung}. Probabilistic carry processes also have an extensive
combinatorial theory \cite{diaconisfulman}. The object here is the finite
response determined by \emph{fixed} operands. The residue and both
outgoing carries are part of its specification.

\begin{theorem}\label{thm:count}
For every integer $B\geq2$,
\[
 |\Q|=5B-4.
\]
For a fixed zero-input residue $v=\res{C+S+H}$, there are five states
when $v<B-2$, and three states at each of $v=B-2$ and $v=B-1$.
All these states have constructive operand representatives.
\end{theorem}

\section{Normal forms and exact classification}
Write
\[
 C+S=qB+r,\qquad r+H=kB+v,
\]
where $0\leq r,v<B$ and $q,k\in\{0,1\}$. Let
$T_{C,S,H}(a,b)=(c_1,c_2)$ be the carry part of the response.

\begin{lemma}\label{lem:residue}
For every incoming pair,
\begin{equation}\label{eq:residue}
 \rho=\res{v+a+b}.
\end{equation}
If $r<B-1$, then
\begin{equation}\label{eq:regular}
 T(a,b)=\left(q,\left\lfloor\frac{r+H+a+b}{B}\right\rfloor\right).
\end{equation}
If $r=B-1$, then $q=0$ and
\begin{equation}\label{eq:special}
 T(a,b)=\left(a,\left\lfloor\frac{[B-1+a]_B+H+b}{B}\right\rfloor\right).
\end{equation}
\end{lemma}
\begin{proof}
Removing the first multiple of $B$ does not change the final residue,
which proves \eqref{eq:residue}. If $r<B-1$, adding $a$ cannot change
$q$. If $r=B-1$, the bound $C+S\leq2B-2$ forces $q=0$, and
$[B-1+a]_B$ equals $B-1$ at $a=0$ and zero at $a=1$.
These observations give \eqref{eq:regular} and \eqref{eq:special}.
\end{proof}

\begin{table}[ht]
\centering
\caption{The complete carry normal forms. Each two-digit entry is the
ordered pair $(c_1,c_2)$. The residue is given by \eqref{eq:residue}.}
\label{tab:forms}
\begin{tabular}{@{}ccrrrr@{}}
\toprule
Residue $v$ & Type $t$ & $T(00)$ & $T(01)$ & $T(10)$ & $T(11)$\\
\midrule
$0\leq v\leq B-3$ &0&00&00&00&00\\
 &1&01&01&01&01\\
 &2&01&01&10&10\\
 &3&10&10&10&10\\
 &4&11&11&11&11\\
\midrule
$B-2$&0&00&00&00&01\\
 &1&01&01&10&11\\
 &2&10&10&10&11\\
\midrule
$B-1$&0&00&01&01&01\\
 &1&00&01&10&10\\
 &2&10&11&11&11\\
\bottomrule
\end{tabular}
\end{table}

\begin{proof}[Proof of Theorem~\ref{thm:count}]
Suppose first that $v\leq B-3$. In the regular case $r<B-1$, neither
$v+1$ nor $v+2$ reaches $B$. Thus $T$ is the constant pair $(q,k)$.
Each of the four pairs occurs: choose $r=v$, $H=0$ for $k=0$;
choose $r=v+1$, $H=B-1$ for $k=1$. Both choices have $r\leq B-2$.
For either choice, $q=0$ is realized by $(C,S)=(0,r)$ and $q=1$
by $(C,S)=(B-1,r+1)$.

The remaining case has $r=B-1$ and $H=v+1$. It gives
$T(a,b)=(a,1-a)$, the fifth map. This case is realized by
$(C,S)=(0,B-1)$. The five maps are distinct.

Now let $v\in\{B-2,B-1\}$. In the regular case,
$r+H\leq2B-3$ excludes $r+H=B+v$. Hence $k=0$.
For $v=B-2$, the second carry is one exactly when $a=b=1$.
For $v=B-1$, it is one exactly when $a+b\geq1$.
Either value of $q$ occurs by taking $r=0$, $H=v$, and
$(C,S)=(0,0)$ or $(B-1,1)$.

The special case $r=B-1$ has $H=B-1$ at $v=B-2$ and $H=0$
at $v=B-1$. It supplies the third map in each terminal group.
Table~\ref{tab:forms} lists every result. Different $v$ values are
separated by the residue at input $(0,0)$; maps within each group are
separated by their carry tables. Therefore the total is
$5(B-2)+3+3=5B-4$. When $B=2$, the middle range is empty and the
six terminal maps still give the stated count.
\end{proof}

\section{Canonical encoding and the exact bit bound}
For the types in Table~\ref{tab:forms}, define
\begin{equation}\label{eq:label}
 \ell(v,t)=
 \begin{cases}
 5v+t,&0\leq v<B-2,\quad 0\leq t<5,\\
 5(B-2)+3(v-B+2)+t,&B-2\leq v<B,\quad 0\leq t<3.
 \end{cases}
\end{equation}
The labels are precisely $0,\ldots,5B-5$. They require no table indexed
by the original operand triple.

An encoder first computes $(q,r)$ and $(k,v)$ above. For $v<B-2$,
type $2$ applies when $r=B-1$. Otherwise $t=k$ if $q=0$ and
$t=3+k$ if $q=1$. At the terminal residues, type $1$ applies when
$r=B-1$; otherwise $t=2q$. A decoder recovers $(v,t)$ by division
with remainder by five or three, reads the corresponding carry row,
and returns \eqref{eq:residue} with that row's carry pair.

\begin{corollary}\label{cor:bits}
The minimum number of bits in a fixed-length code that recovers the
complete response is
\[
 b(B)=\left\lceil\log_2(5B-4)\right\rceil.
\]
For $B=2^w$, this is $3$ at $w=1$, $4$ at $w=2$, and $w+3$ for
all $w\geq3$.
\end{corollary}
\begin{proof}
A code of length $b$ has at most $2^b$ values. Distinct functions in
$\Q$ must have different codes, giving the lower bound.
The labels in \eqref{eq:label} achieve it. For $B>4$,
$4B<5B-4<8B$, which gives $w+3$ in the binary case.
\end{proof}

A byte stage has $1276$ states and needs eleven bits. A 32-bit word
has $21\,474\,836\,476$ states and needs exactly 35 bits.
The input consists of 96 operand bits, but the preserved object is the
four-input response, not the three operands themselves.
This is a representation bound; encoder and decoder execution costs are
separate from the number of state bits.

For example, at $B=256$, the operands $(37,218,13)$ have $r=255$
and $v=12$. Their label is $5\cdot12+2=62$. In input order
$00,01,10,11$, the full outputs are
\[
 (12,0,1),\quad(13,0,1),\quad(13,1,0),\quad(14,1,0).
\]
Exactly 256 operand triples produce this same function.

\section{Class multiplicities and carry images}
The quotient states are not equally likely under uniform operands.
Their exact multiplicities follow from the triangular distribution
of $C+S$.

\begin{proposition}\label{prop:weights}
For $0\leq v\leq B-3$, the numbers of operand triples of types
$0,1,2,3,4$ are respectively
\begin{align*}
 n_0(v)&=\frac{(v+1)(v+2)}2,&
 n_1(v)&=\frac{B(B-1)-(v+1)(v+2)}2,\\
 n_2(v)&=B,&
 n_3(v)&=\frac{(v+1)(2B-2-v)}2,\\
 n_4(v)&=\frac{(B-v-2)(B-v-1)}2.
\end{align*}
At either terminal residue, the three multiplicities are
$B(B-1)/2$, $B$, and $B(B-1)/2$.
\end{proposition}
\begin{proof}
For $r\leq B-2$, the number of pairs with $C+S=r$ is $r+1$;
the number with $C+S=B+r$ is $B-r-1$. Once $r$ and $v$ are fixed,
$H=[v-r]_B$ is unique. For $k=0$, sum the appropriate pair count
from $r=0$ to $r=v$. For $k=1$, sum from $r=v+1$ to $r=B-2$.
These four sums give $n_0,n_3,n_1,n_4$. The special case
$C+S=B-1$ has exactly $B$ pairs. At a terminal residue only $k=0$
occurs in the regular case, so both sums extend through $r=B-2$.
\end{proof}

Let $\kappa(C,S,H)=|\im T_{C,S,H}|$ be the number of distinct outgoing
carry pairs. This cardinality concerns the carry projection only.
The complete responses also retain the residues.

\begin{corollary}\label{cor:rank}
Among all $B^3$ operand triples, the counts for $\kappa=1,2,3$ are
\[
 B(B-1)(B-2),\qquad 3B^2-4B,\qquad 2B,
\]
and $\kappa=4$ does not occur. Under independent uniform operands,
\begin{equation}\label{eq:prob}
 \Pr(\kappa>1)=\frac3B-\frac2{B^2},\qquad
 \mathbb E\kappa=1+\frac3B.
\end{equation}
\end{corollary}
\begin{proof}
The four regular middle types have constant carry maps. Their total
multiplicity at each middle residue is $B(B-1)$. The middle special
type has image size two and contributes $B(B-2)$ triples in total.
The four regular terminal types contribute $2B(B-1)$ more triples of
image size two. Each special terminal type has image size three and
multiplicity $B$. The stated counts follow. Dividing their sums by
$B^3$ gives \eqref{eq:prob}.
\end{proof}

For $B=256$, these counts are $16\,581\,120$, $195\,584$, and $512$.
For binary words, \eqref{eq:prob} becomes
$\Pr(\kappa>1)=3\cdot2^{-w}-2\cdot4^{-w}$ and
$\mathbb E\kappa=1+3\cdot2^{-w}$.
These probabilities use uniform \emph{operand triples}, rather than
uniform labels or independent samples from the normal-form table.

\section{Independently recoverable digit stages}
The whole-word response and a collection of separately queryable digit
responses specify different objects. The distinction determines the
encoding bound.

Take radices $B_1,\ldots,B_m\geq2$, from least to most significant.
At stage $i$, select an independent operand triple of radix $B_i$ and
retain its complete function in $\mathcal Q_{B_i}$. The contract permits
any of the four carry inputs to be queried at \emph{each} stage.

\begin{theorem}\label{thm:stages}
The number of independently recoverable stage tuples is exactly
\[
 \prod_{i=1}^{m}(5B_i-4).
\]
Their optimal fixed-length encoding uses
$\left\lceil\log_2\prod_i(5B_i-4)\right\rceil$ bits.
The tuple determines the full composed residue and both outgoing
carries at every intermediate boundary, for any initial carry pair.
\end{theorem}
\begin{proof}
Every stage function has a representative by
Theorem~\ref{thm:count}. Independent digit choices therefore realize
every tuple in the Cartesian product. Distinct tuples differ at a
stage that the contract allows a recipient to query, so they require
distinct labels. Mixed-radix encoding of the stage labels attains the
bound. To compose them, evaluate the least significant stage and feed
its outgoing pair into the next stage. Repeat and assemble the residue
digits with weights $1,B_1,B_1B_2,\ldots$.
\end{proof}

Four byte stages have $1276^4=2\,650\,957\,086\,976$ tuples and
require 42 bits. Four separate eleven-bit fields use 44 bits.
A single 32-bit whole-word response requires 35 bits instead.
The 42-bit minimum is for independently recoverable stage functions.
It is not a lower bound for an interface that exposes only trajectories
reachable from the four initial whole-word carry pairs.

\section{Reproduction and scope}
The supplied Python encoder implements \eqref{eq:label} for arbitrary
integer radices. Its decoder uses Table~\ref{tab:forms}; a separate
forward oracle evaluates \eqref{eq:first}--\eqref{eq:second} directly.
The verification compares all four outputs for every operand triple
at radices $2$ through $20$: $44\,099$ triples and $176\,396$ responses.
It also checks every class size and constructive representative.

A separately compiled C++ program exhausts all binary widths one
through eight. At width eight it compares all $16\,777\,216$ triples
and obtains the class and carry-image counts above. Further exact
checks cover $1340$ operand triples through 512-bit words and $2560$
four-input stage-composition comparisons, including nonbinary mixed
radices. All comparisons pass. The complete-function experiments
support the implementations; the case classification proves the
statements for every radix.

The figure uses the complete radix-eight example: 512 operand triples
map to 36 response states. Every link comes from the encoder, and each
state has its actual multiplicity. The manuscript and source package
state the response contract explicitly. The encoded state preserves
that contract rather than the original operands, arithmetic timing,
or carries inside an unpartitioned word.

\begin{thebibliography}{9}
\bibitem{brentkung} R. P. Brent and H. T. Kung,
A regular layout for parallel adders,
\emph{IEEE Transactions on Computers} C-31 (1982), 260--264.
\href{https://doi.org/10.1109/TC.1982.1675982}{doi:10.1109/TC.1982.1675982}.
\bibitem{diaconisfulman} P. Diaconis and J. Fulman,
Carries, shuffling, and symmetric functions,
\emph{Advances in Applied Mathematics} 43 (2009), 176--196.
\href{https://doi.org/10.1016/j.aam.2009.02.002}{doi:10.1016/j.aam.2009.02.002}.
\end{thebibliography}
\end{document}
